% GENERATED FILE. Edit canonical lessons, then run npm run generate:books.
% canonical-source-hash: fnv1a64:8ecd23d5f705df8f
% canonical-lessons: GU-C492-piyar, GU-C492-sasru, GU-C492-bhanej, GU-C492-pautra, GU-C492-pitrai

\chapter{Two Homes}
\label{ch:gu-two-homes}
\hlchaptermodality{492}

\begin{chapteropening}
\textbf{By the end of this chapter:} \emph{I can talk about a married woman's two homes.}

Complete the last lesson of chapter 492: I can talk about a married woman's two homes.
\end{chapteropening}

\section[\texorpdfstring{\gu{પિયર} (piyar)}{પિયર (piyar)}]{\gu{પિયર} (piyar) — a married woman's parents' home}

\label{lesson:GU-C492-piyar}



\begin{quote}
Before the new one: say the Gujarati for the father-in-law of one's child, then the Gujarati for the mother-in-law of one's child.
\end{quote}

\subsection*{You'll want to know: \gu{પિયર}}
\textbf{\gu{પિયર}} — \emph{piyar} — \textquotedblleft{}a married woman's parents' home\textquotedblright{}.

\begin{quote}
\gu{દિવાળીમાં} \gu{મીરા} \gu{પિયર} \gu{જશે}.
\end{quote}

\begin{quote}
\emph{divāḷīmā̃ Mīrā piyar jaśe.}
\end{quote}

\begin{quote}
At Diwali Mira will go to her parents' home.
\end{quote}

A neuter noun.

\begin{cousinweb}
\textbf{\gu{પિયર}} (\emph{piyar}) goes back to a reconstructed \emph{*pitṛghara-}, \textquotedblleft{}father's house\textquotedblright{}; CDIAL lists the Gujarati word (CDIAL 8180).
\end{cousinweb}

\begin{grammarlens}[title={What you've built}]
One more word for describing what you lived and felt.
\end{grammarlens}

\subsection*{Your turn}
\begin{itemize}
\raggedright
  \item \emph{Say it:} \emph{piyar}
  \item \emph{Say it:} \emph{piyar}, once more
  \item \emph{Say it:} \emph{divāḷīmā̃ Mīrā piyar jaśe.} — At Diwali Mira will go to her parents' home.
  \item \emph{Recall:} say the Gujarati for the father-in-law of one's child, then the Gujarati for the mother-in-law of one's child, then say \emph{piyar} again
\end{itemize}

\begin{tcolorbox}[breakable,colback=teal!4,colframe=teal!35!black,title={Before you move on}]
What does \gu{પિયર} mean? (\textquotedblleft{}A married woman's parents' home\textquotedblright{}.) Say it once more.
\end{tcolorbox}
\section[\texorpdfstring{\gu{સાસરું} (sāsrũ)}{સાસરું (sāsrũ)}]{\gu{સાસરું} (sāsrũ) — the in-laws' home}

\label{lesson:GU-C492-sasru}



\begin{quote}
Before the new one: say the Gujarati for the mother-in-law of one's child, then the Gujarati for a married woman's parents' home.
\end{quote}

\subsection*{You'll want to know: \gu{સાસરું}}
\textbf{\gu{સાસરું}} — \emph{sāsrũ} — \textquotedblleft{}the in-laws' home\textquotedblright{}.

\begin{quote}
\gu{લગ્ન} \gu{પછી} \gu{દીકરી} \gu{સાસરે} \gu{ગઈ}.
\end{quote}

\begin{quote}
\emph{lagna pachī dīkrī sāsre gaī.}
\end{quote}

\begin{quote}
After the wedding the daughter went to her in-laws' home.
\end{quote}

A neuter noun.

\begin{cousinweb}
\textbf{\gu{સાસરું}} (\emph{sāsrũ}) is inherited from Sanskrit \emph{śvāśura}, \textquotedblleft{}belonging to the father-in-law\textquotedblright{}; CDIAL lists the Gujarati word, \textquotedblleft{}father-in-law's house\textquotedblright{} (CDIAL 12767).
\end{cousinweb}

\begin{grammarlens}[title={What you've built}]
One more word for describing what you lived and felt.
\end{grammarlens}

\subsection*{Your turn}
\begin{itemize}
\raggedright
  \item \emph{Say it:} \emph{sāsrũ}
  \item \emph{Say it:} \emph{sāsrũ}, once more
  \item \emph{Say it:} \emph{lagna pachī dīkrī sāsre gaī.} — After the wedding the daughter went to her in-laws' home.
  \item \emph{Recall:} say the Gujarati for the mother-in-law of one's child, then the Gujarati for a married woman's parents' home, then say \emph{sāsrũ} again
\end{itemize}

\begin{tcolorbox}[breakable,colback=teal!4,colframe=teal!35!black,title={Before you move on}]
What does \gu{સાસરું} mean? (\textquotedblleft{}The in-laws' home\textquotedblright{}.) Say it once more.
\end{tcolorbox}
\section[\texorpdfstring{\gu{ભાણેજ} (bhāṇej)}{ભાણેજ (bhāṇej)}]{\gu{ભાણેજ} (bhāṇej) — a sister's child}

\label{lesson:GU-C492-bhanej}



\begin{quote}
Before the new one: say the Gujarati for a married woman's parents' home, then the Gujarati for the in-laws' home.
\end{quote}

\subsection*{You'll want to know: \gu{ભાણેજ}}
\textbf{\gu{ભાણેજ}} — \emph{bhāṇej} — \textquotedblleft{}a sister's child\textquotedblright{}.

\begin{quote}
\gu{મામા} \gu{ભાણેજ} \gu{માટે} \gu{રમકડું} \gu{લાવ્યા}.
\end{quote}

\begin{quote}
\emph{māmā bhāṇej māṭe ramakḍũ lāvyā.}
\end{quote}

\begin{quote}
Uncle brought a toy for his sister's child.
\end{quote}

A masculine or feminine noun, by the person meant.

\begin{cousinweb}
\textbf{\gu{ભાણેજ}} (\emph{bhāṇej}) is inherited from Sanskrit \emph{bhāginēya}, \textquotedblleft{}sister's son\textquotedblright{}; CDIAL lists the Gujarati word (CDIAL 9433).
\end{cousinweb}

\begin{grammarlens}[title={What you've built}]
One more word for describing what you lived and felt.
\end{grammarlens}

\subsection*{Your turn}
\begin{itemize}
\raggedright
  \item \emph{Say it:} \emph{bhāṇej}
  \item \emph{Say it:} \emph{bhāṇej}, once more
  \item \emph{Say it:} \emph{māmā bhāṇej māṭe ramakḍũ lāvyā.} — Uncle brought a toy for his sister's child.
  \item \emph{Recall:} say the Gujarati for a married woman's parents' home, then the Gujarati for the in-laws' home, then say \emph{bhāṇej} again
\end{itemize}

\begin{tcolorbox}[breakable,colback=teal!4,colframe=teal!35!black,title={Before you move on}]
What does \gu{ભાણેજ} mean? (\textquotedblleft{}A sister's child\textquotedblright{}.) Say it once more.
\end{tcolorbox}
\section[\texorpdfstring{\gu{પૌત્ર} (pautra)}{પૌત્ર (pautra)}]{\gu{પૌત્ર} (pautra) — a grandson (a son's son)}

\label{lesson:GU-C492-pautra}



\begin{quote}
Before the new one: say the Gujarati for the in-laws' home, then the Gujarati for a sister's child.
\end{quote}

\subsection*{You'll want to know: \gu{પૌત્ર}}
\textbf{\gu{પૌત્ર}} — \emph{pautra} — \textquotedblleft{}a grandson (a son's son)\textquotedblright{}.

\begin{quote}
\gu{દાદાને} \gu{પૌત્ર} \gu{સાથે} \gu{રમવું} \gu{ગમે} \gu{છે}.
\end{quote}

\begin{quote}
\emph{dādāne pautra sāthe ramvũ game che.}
\end{quote}

\begin{quote}
Grandfather likes playing with his grandson.
\end{quote}

A masculine noun.

\begin{cousinweb}
\textbf{\gu{પૌત્ર}} (\emph{pautra}) is a tatsama, Sanskrit \emph{pautra}, \textquotedblleft{}a son's son, a grandson\textquotedblright{}, as Monier-Williams records it.
\end{cousinweb}

\begin{grammarlens}[title={What you've built}]
One more word for describing what you lived and felt.
\end{grammarlens}

\subsection*{Your turn}
\begin{itemize}
\raggedright
  \item \emph{Say it:} \emph{pautra}
  \item \emph{Say it:} \emph{pautra}, once more
  \item \emph{Say it:} \emph{dādāne pautra sāthe ramvũ game che.} — Grandfather likes playing with his grandson.
  \item \emph{Recall:} say the Gujarati for the in-laws' home, then the Gujarati for a sister's child, then say \emph{pautra} again
\end{itemize}

\begin{tcolorbox}[breakable,colback=teal!4,colframe=teal!35!black,title={Before you move on}]
What does \gu{પૌત્ર} mean? (\textquotedblleft{}A grandson (a son's son)\textquotedblright{}.) Say it once more.
\end{tcolorbox}
\section[\texorpdfstring{\gu{પિતરાઈ} (pitrāī)}{પિતરાઈ (pitrāī)}]{\gu{પિતરાઈ} (pitrāī) — a cousin (on the father's side)}

\label{lesson:GU-C492-pitrai}



\begin{quote}
Before the new one: say the Gujarati for a sister's child, then the Gujarati for a grandson.
\end{quote}

\subsection*{You'll want to know: \gu{પિતરાઈ}}
\textbf{\gu{પિતરાઈ}} — \emph{pitrāī} — \textquotedblleft{}a cousin (on the father's side)\textquotedblright{}.

\begin{quote}
\gu{મારો} \gu{પિતરાઈ} \gu{ભાઈ} \gu{મારી} \gu{સાથે} \gu{ભણે} \gu{છે}.
\end{quote}

\begin{quote}
\emph{māro pitrāī bhāī mārī sāthe bhaṇe che.}
\end{quote}

\begin{quote}
My cousin studies with me.
\end{quote}

A masculine or feminine noun, by the person meant.

\begin{cousinweb}
\textbf{\gu{પિતરાઈ}} (\emph{pitrāī}) goes back to a reconstructed \emph{*pitryajāta-}, \textquotedblleft{}father's brother's son\textquotedblright{}; CDIAL lists the Gujarati word, \textquotedblleft{}relating to father's brother\textquotedblright{} (CDIAL 8193).
\end{cousinweb}

\begin{grammarlens}[title={What you've built}]
One more word for describing what you lived and felt.
\end{grammarlens}

\subsection*{Your turn}
\begin{itemize}
\raggedright
  \item \emph{Say it:} \emph{pitrāī}
  \item \emph{Say it:} \emph{pitrāī}, once more
  \item \emph{Say it:} \emph{māro pitrāī bhāī mārī sāthe bhaṇe che.} — My cousin studies with me.
  \item \emph{Recall:} say the Gujarati for a sister's child, then the Gujarati for a grandson, then say \emph{pitrāī} again
\end{itemize}

\begin{tcolorbox}[breakable,colback=teal!4,colframe=teal!35!black,title={Before you move on}]
What does \gu{પિતરાઈ} mean? (\textquotedblleft{}A cousin (on the father's side)\textquotedblright{}.) Say it once more.
\end{tcolorbox}
